\documentclass{article}
\usepackage[utf8]{inputenc}
\usepackage[T1]{fontenc}
\usepackage{geometry}
\geometry{
    a4paper,
    left=2.5cm,
    right=2.5cm,
    top=2.5cm,
    bottom=2.5cm
}
\usepackage{graphicx}
\usepackage{float}
\usepackage{hyperref}
\usepackage[english]{babel}
\usepackage{amsmath}
\usepackage{amssymb}
\usepackage{array}
\usepackage{booktabs}
\usepackage{caption}
\usepackage{subcaption}
\usepackage{enumitem}
\usepackage{fancyhdr}
\usepackage{titlesec}
\usepackage{url}
\usepackage{textcomp}
\usepackage{xcolor}
\usepackage{listings}
\usepackage{multirow}
\usepackage{longtable}
\usepackage{tabularx}
\usepackage{times}

% Set up page style
\pagestyle{fancy}
\fancyhf{}
\rfoot{\thepage}
\renewcommand{\headrulewidth}{0pt}

% Title formatting
\titleformat{\chapter}[display]{\normalfont\bfseries}{\chaptertitle\thechapter}{0pt}{\Huge}
\titlespacing*{\chapter}{0pt}{50pt}{40pt}

% Document info
\title{Template Document}
\author{Academic Support Agent}
\date{\today}

\begin{document}
\maketitle

\addtocontents{toc}{\protect\thispagestyle{empty}}
%IN CIO DA CAPA DA MONOGRAFIA
\thispagestyle{empty}
\begin{center}
\textbf{ UNIVERSIDADE DO ESTADO DO AMAZONAS \\
ESCOLA SUPERIOR DE TECNOLOGIA \\[50pt] }
\textbf{ \\[70pt] {\bigsize CARLA PATRICIA MICHILES ARAUJO} \\[120pt] }
\textbf{ {\bigsize SISTEMA DE AQUISI O DE SINAIS ELETROCARDIOGR FICOS UTILIZANDO TECNOLOGIA WIRELESS}  \\[104pt] }
\end{center}
% COMENT RIOS DA FOLHA DE ROSTO
\hspace*{8cm}

\vspace*{\fill}

\begin{center}
Manaus \\ 2017
\end{center}

%FIM DA CAPA DA MONOGRAFIA

%IN CIO DA FOLHA DE ROSTO
\newpage
\thispagestyle{empty}
\begin{center}


% \textbf{ \\[70pt] {\bigsize CARLA PATRICIA MICHILES ARAUJO} \\[120pt] }
\textbf{ {\bigsize CARLA PATRICIA MICHILES ARAUJO} \\[120pt] }
\textbf{ {\bigsize SISTEMA DE AQUISI O DE SINAIS ELETROCARDIOGR FICOS UTILIZANDO TECNOLOGIA WIRELESS}  \\[50pt] }
\end{center}
% COMENT RIOS DA FOLHA DE ROSTO
\hspace*{8cm}
\begin{flushright}
\begin{minipage}{8cm}
\begin{singlespace}
Projeto de pesquisa desenvolvido durante a disciplina de Trabalho de Conclus o de Curso II e apresentada   banca avaliadora do Curso de Engenharia El trica da Escola Superior de Tecnologia da Universidade do Estado do Amazonas, como pr -requisito para obten o do t tulo de Engenheiro Eletricista.\\[50pt]
\end{singlespace}
\end{minipage}
\end{flushright}
\begin{center}
Orientador: Wheidima Carneiro de Melo, Me.

\end{center}

\vspace*{\fill}
\begin{center}
Manaus \\ 2017
\end{center}



%FIM DA FOLHA DE ROSTO

%INICIO DA PAG CATALOGA O

\newpage
\thispagestyle{empty}

%\begin{figure}[!htb]
%\begin{center}
%			\includegraphics[width=\textwidth]{catt.PNG}
%            \vspace*{\fill} 
%            \begin{quote} 
%            \centering  
%            \end{quote}
%            \vspace*{\fill}
%			\label{fig:cat.PNG}
%\end{center}
%\end{figure}

\includepdf[pages=-]{carlaze.pdf}

\newpage
\thispagestyle{empty}
\begin{center}


% \textbf{ \\[30pt] {\bigsize CARLA PATRICIA MICHILES ARAUJO} \\[50pt] }
\textbf{ {\bigsize CARLA PATRICIA MICHILES ARAUJO} \\[50pt] }
\textbf{ {\bigsize SISTEMA DE AQUISI O DE SINAIS ELETROCARDIOGR FICOS UTILIZANDO TECNOLOGIA WIRELESS}  \\[50pt] }
\end{center}
% COMENT RIOS DA FOLHA DE ROSTO
\hspace*{8cm}
\begin{flushright}
\begin{minipage}{8cm}
\begin{singlespace}
Pesquisa desenvolvida durante a disciplina de Trabalho de Conclus o de Curso II e apresentada   banca avaliadora do Curso de Engenharia El trica da Escola Superior de Tecnologia da Universidade do Estado do Amazonas, como pr -requisito para obten o do t tulo de Engenharia Eletricista.\\[30pt]
\end{singlespace}
\end{minipage}
\end{flushright}
\begin{center}
Nota obtida:\uline{\hspace{1cm}}(\uline{\hspace{8cm}})
\newline \\
Aprovada em 22/11/2017
\newline\\
 rea de concentra o: Engenharia Biom dica \,\,\,\,\,\,\,\,\,\,\,\,\,\,\,\,\,\,\,
\end{center}
\vspace{1.2cm}
\begin{center}
BANCA EXAMINADORA
\end{center}
\vspace{1.0cm}

%\begin{flushright}
\begin{minipage}[l]{14cm}
\begin{center}
\hspace{1cm}\uline{\hspace{10.5cm}} \\
\hspace{1cm}Orientador: Wheidima Carneiro de Melo, Me.

\hspace{1cm}\uline{\hspace{10.5cm}} \\
\hspace{1cm}Avaliador: Jozias Parente de Oliveira, Dr.

\hspace{1cm}\uline{\hspace{10.5cm}} \\
\hspace{1cm}Avaliador: Charles Luiz de Melo, Me.

\end{center}
\end{minipage}

%\end{flushright}
\hspace*{8cm}

\vspace*{\fill}
\begin{center}
Manaus\\2017
\end{center}



\newpage
\thispagestyle{empty}
\begin{flushright}
\begin{minipage}{8cm}
\begin{singlespace}
{\begin{center}\vspace{18cm}\textbf{\normalsize Dedicat ria}\vspace{36pt}\end{center}}
Dedico este trabalho ao meu irm o Pablo (in memorian) que est  presente a toda hora no meu cora o e aos momentos de alegria que me proporcionou na sua caminhada.\\[50pt]
\end{singlespace}
\end{minipage}
\end{flushright}
\newpage
\thispagestyle{empty}
\vspace*{4.5cm}
{\begin{center}\textbf{\normalsize AGRADECIMENTO}\vspace{36pt}\end{center}}
\hspace*{0.8cm}  Deus, que tornou poss vel, mesmo com muitos sacrif cios me deu f  para finaliza o deste trabalho.

\hspace*{0.8cm} Agrade o   minha fam lia, que sempre proporcionaram condi es adequadas para crescer na vida pessoal e profissional e ao meu noivo, amigo e compaheiro de todas as horas, Kenny Vinente, pelo carinho, paci ncia, compreens o e amor.

\hspace*{0.8cm} Agrade o ao meu orientador, Wheidima Carneiro de Melo, por fornecer todo o apoio que eu precisava para a realiza o deste trabalho.


\hspace*{0.8cm}Aos engenheiros, Thales Araujo, Ronald Santos e Adriano Ibrahim, pelo tempo dedicado na empresa Philips, que contribuiram com material e suporte no entendimento te rico do sinal do eletrocardiograma.


\hspace*{0.8cm}Aos alunos de gradua o, Giovanni Antonnacio e Marcus Vin cius pelo suporte no desenvolvimento do projeto e aos Not veis que estiveram ao meu lado me incentivando durante todo o processo.




\newpage


\thispagestyle{empty}
\vspace*{4.4cm}
{\begin{center}\textbf{\normalsize RESUMO}\vspace{36pt}\end{center}}
\textcolor{red}{Mude isso! No resumo vc dar  uma vis o geral do trabalho. Ent o   como se vc colocasse um pouquinho da intro, desen e conclu...sacou?}
O presente trabalho tem como objetivo desenvolver um prot tipo de eletrocardi grafo port til de baixo custo que realize a aquisi o e transmiss o dos sinais ECG em tempo real. Sua apresenta o est  dividida em quatro cap tulos. O primeiro cap tulo, Referencial Te rico, apresenta assuntos referente a composi o do eletrocardiograma, necess rias para o desenvolvimento deste trabalho, o cora o, o ciclo card aco, o eletrocardiograma, deriva es eletrocardiogr ficas, o tri ngulo de einthoven, eletrodos e interfer ncia em sinais de ECG. O segundo cap tulo, Metodologia, apresenta os materias e m todos necess rio para aquisi o e transmiss o dos sinais de ECG. O terceiro cap tulo, Implementa o do Projeto, descreve o uso dos componentes, assim como o desenvolvimento do hardware utilizando a placa de aquisi o e o Rapsberry Pi 3, al m de descrever o uso dos protocolos utilizados, o funcionamento da arquitetura cliente-servidor para transmiss o de dados utilizando a rede Wi-Fi como meio de comunica o. O quarto cap tulo apresenta os resultados obtidos durante o desenvolvimento do projeto. Apesar do estudo ter um car ter de implementa o do circuito de aquisi o ECG, os resultados obtidos apresentaram ru dos nas formas de ondas, que poderiam ser diminu dos se algoritmos espec ficos com certo grau de complexidade fossem utilizados para realizar a separa o do sinal e ru do. Nas conclus es s o apresentadas as dificuldades enfrentadas na realiza o deste projeto, os objetivos que foram alcan ados e aplica es para o prot tipo desenvolvido.


\textbf{Palavras chave}: eletrocardiograma. prot tipo de baixo custo, Raspberry Pi 3 \textcolor{red}{Mude tbm! Aqui tem q ser algo que esteja diretamente ligado ao seu trabalho. SUGEST O: SISTEMA DE AQUISI O DE ECG, DISPOSITIVO PORT TIL, ENGENHARIA BIOM DICA}.  
\newpage
\thispagestyle{empty}
\vspace*{4.4cm}
{\begin{center}\textbf{\normalsize ABSTRACT}\vspace{36pt}\end{center}}
The present work aims to develop a prototype of a low cost portable electrocardiograph that performs the acquisition and transmission of ECG signals in real time. Its presentation is divided into four chapters. The first chapter, Theoretical Framework, presents subjects related to the composition of the electrocardiogram necessary for the development of this work, the heart, the cardiac cycle, the electrocardiogram, electrocardiographic leads, the eithoven s triangle, electrodes and interference on ECG signals. The second chapter, Methodology, presents the materials and methods necessary for the acquisition and transmission of ECG signals. The third chapter, Project Implementation, describes the use of the components, as well as hardware development using the acquisition board and a Rapsberry Pi 3. In addition, it is described the protocols used and the operation of the client-server architecture for data transmission using the wireless network to communicate them. The fourth chapter presents the results obtained during the development of the project. Although the study focused on the implementation of an ECG acquisition circuit, the results obtained presented noises in the waveforms, which could be reduced if specific algorithms with a certain degree of complexity were used to perform signal and noise separation. Finally, the conclusion presents the challenges faced in the accomplishment of this project, the objectives achieved and applications for the developed prototype.


\textbf{Keywords}: electrocardiogram, low cost prototype, Rapsberry Pi 3


%IN CIO DA LISTA DE FIGURAS

\newpage
\pagestyle{empty}
\vspace*{4cm}


%\pagestyle{plain}
%FIM DA LISTA DE FIGURAS



%IN CIO DO SUM RIO
\newpage
\pagestyle{empty}
\vspace*{3.5cm}
\renewcommand{\contentsname}{\begin{center}\textbf{\normalsize SUM RIO}\vspace{36pt}\end{center}}


%FIM DO SUM RIO
\pagestyle{myheadings}

%IN CIO DO CAP TULO 1 - INTRODU O
\newpage
\vspace*{4cm}

%\fancypagestyle{plain}{
%  \fancyhf{}% Clear header/footer
%  \fancyhead[R]{}% Right header
%  \fancyfoot[L]{}% Left footer
%  \fancyfoot[R]{\thepage}% Right footer
%}

%\thispagestyle{fancy}

%\pagestyle{fancy}
%\renewcommand{\headrulewidth}{0pt}
%\rhead{\thepage}
%\fancyfoot[C]{}

%\fancyfoot[R]{\thepage}
\setcounter{page}{12}
\begin{center}
\textbf{INTRODU O\\}
\end{center}
\par
\addcontentsline{toc}{section}{INTRODU \~AO}	

	Os diversos processos fisiol gicos no corpo humano s o respons veis por gerar sinais biol gicos. Tais sinais podem ser el tricos, qu micos ou ac sticos em suas origens, e a an lise dos mesmos   proveitosa na identifica o, tratamento e entendimento de condi es patol gicas do corpo humano \textcolor{red}{aqui vai uma refer ncia}. Entre os principais sinais biol gicos, pode-se destacar o sinal de eletrocardiograma (ECG), que registra a varia o dos est mulos el tricos aplicados aos  trios e ventr culos do cora o, devido a sua ampla difus o, progn sticos bem estabelecidos e t cnica n o invasiva \textcolor{red}{aqui vai uma refer ncia}.

	A import ncia em realizar estudos no sinal de ECG pode ser explicada devido   quantidade de informa es que podem ser extra das desse sinal. Estas caracter sticas s o utilizadas pelo especialista com o objetivo de identificar se o indiv duo tem alguma doen a card aca ou at  mesmo outros tipos de patologias n o card acas. A an lise do sinal de ECG pode ser feita diretamente pela sua forma de onda, portanto,   de suma import ncia a boa capta o deste sinal realizada por equipamentos espec ficos. 

	Entretanto, apesar da larga utiliza o de eletrocardi grafos (equipamento que adquire o sinal de ECG), poucos destes empregam tecnologia Wireless em seus sistemas de monitora o de pacientes \textcolor{red}{aqui vai uma ref}. Levantamentos iniciais na literatura apresentam como motivos o custo elevado do emprego desta tecnologia e tamb m a necessidade de transmiss o eficiente dos sinais, sem interfer ncia de ru dos inseridos pelo meio de comunica o \textcolor{red}{aqui vai uma ref}. Al m disso, observou-se que historicamente estes tipos de dispositivos n o s o do tipo port til, possibilitando um mercado promissor. 
    
Com o avan o da tecnologia existente no mercado,   poss vel adquirir amplificadores de instrumenta o, filtros, microcontroladores, conversores anal gico-digital e utiliz -los para a manipula o de sinais, sendo uma poss vel aplica o comunicar sinais de eletrocardiograma a partir de sistemas microcontrolados, com gera o do sinal e transmiss o em tempo real, utilizando o meio de comunica o \textcolor{red}{melhor colocar MEIO DE TRANSMISS O} a rede Wi-fi \textcolor{red}{referencia para WI-FI}.
  
	O presente trabalho tem como objetivo o estudo de um sistema de aquisi o de sinais eletrocardiogr ficos, de baixo custo, utilizando tecnologia Wireless, em espec fico, o protocolo 802.11, popularmente conhecido como Wi-Fi, \textcolor{red}{melhore essse texto:} ap s esta etapa consistir  em montar um prot tipo de teste de um eletrocardi grafo, no laborat rio de sistemas embarcados- LSE, localizado no HUB - Tecnologia Inova o.
    
    \textcolor{red}{coloque a justificativa academica e a de pesquisa} Academicamente, o projeto de pesquisa contribui para o acervo na  rea de engenharia biom dica incluindo assuntos relacionados   eletr nica anal gica, eletr nica digital, sinais e sistemas microprocessados e linguagem de programa o.
    
    Esta monografica est  organizada nas seguintes partes:
    
    Cap tulo 1: Referencial Te rico. Nesta se o apresenta o material te rico utilizados durante a pesquisa.
    
    Cap tulo 2: Metodologia. Esta se o est  contida em detalhes os materiais e m todos utilizando durante a pesquisa.
    
    Cap tulo 3: Desenvolvimento do prot tipo. Este cap tulo detalha o processo de implementa o do hardware e software do prot tipo do eletrocardiogr fico.
    
    Cap tulo 4: Resultados obtidos. Neste cap tulo est o os dados e par metros obtidos a partir do prot tipo e an lises dos resultados.
    
   Conclus o: Relata as dificuldades encontradas no decorrer do trabalho. Aborda os resultados obtidos pelo prot tipo desenvolvido no projeto. Finalizando com sugest es para trabalhos futuros com o objetivo do aperfei oamento da proposta.
   
    
    


\newpage
\section{REFERENCIAL TE RICO}	

%\pagestyle{fancy}
%\fancyhead{}
%\renewcommand{\headrulewidth}{0pt}
%\rhead{\thepage}
%\fancyfoot[C]{}

\subsection{O CORA O}
\hspace*{0.8cm}O cora o   o  rg o humano mais importante do sistema circulat rio, pois atua como uma bomba que impulsona o sangue para as demais parte do corpo, o lado direito bombeia sangue para os pulm es, e o esquerdo, bombeia para os org es perif ricos. Sua anatomia   dividida em dois lados  espelhados , esquerdo e direito, os quais suportam diferentes sistemas circulat rios, que bombeiam na mesma sintonia e r tmica \cite{guyton}. Seu sistema   dividido em quatro cavidades, sendo constitu do por dois  trios (direito e esquerdo) e dois ventr culos (direito e esquerdo).   Os  trios s o respons veis pelo armazenamento do sangue enquanto os ventr culos t m como fun o realizar o bombeamento do sangue  \cite{silva}.
A Figura \ref{fig:ramcor} apresenta a estrutura interna do cora o, tais como ramifica es, posicionamento dos atrios e ventr culos. 


\begin{figure}[!htb]
\begin{center}
			\caption{Ramifica o do Cora o.}
			\includegraphics[width=.5\textwidth]{Figuras/cor.png}
            \vspace*{\fill} 
            \begin{quote} 
            \centering 
            Fonte: \cite{ronald} 
            \end{quote}
            \vspace*{\fill}
			\label{fig:ramcor}
\end{center}
\end{figure}

O sangue passa do  trio direito para o ventr culo direito por meio da valva tric spide e do lado do  trio esquerdo para o ventr culo por meio da valva mitral ou conhecida como valva bic spide. O per odo de contra o ou bombeamento ventricular   denominado s stole enquanto a fase de relaxamento   chamada di stole. Esse sistema de contra es r tmicas   precedido por eventos el tricos coordenados e intr nsecos ao cora o, respons veis em gerar os potenciais el tricos do eletrocardiograma, que s o produzidos por c lulas eletricamente excit veis.

A ativa o el trica inicia num ponto do  trio direito denominado de n dulo sinoatrial (AS). O n dulo sinoatrial atua de forma independente regulando a frequ ncia card aca. Entretanto, ele   influenciado por sinais neurol gicos e hormonais, interferindo na frequ ncia de batimento card aca \cite{silva}.

As fibras do n dulo sinusal ou sinoatrial possui o di metro entre 3 a 5 micr metros enquanto as fibras atriais musculares que as contornem possui de 10 a 15 micr metros de di metro. Essas fibras est o conectadas de modo direto, ou seja, qualquer potencial de a o que se inicie no n dulo se difunde imediatamente para a parede do m sculo atrial \cite{guyton}.

\subsection{O CICLO CARDI CO}

\hspace*{0.8cm}O ciclo card aco corresponde ao intervalo do in cio de um batimento card aco at  o in cio do pr ximo batimento. O respons vel por fornecer o ciclo card aco   o nodo sinoatrial localizado na parede lateral superior do  trio direito, pr ximo a abertura da veia superior que gera o potencial de a o que se propaga imediatamente pelo os  trios e aos ventr culos atrav s do feixe A-V com dura o acima de $0,1$s. Em fun o desse retardo os  trios se contraiam antes dos ventr culos, portanto, os  trios bombeiam o sangue para os ventr culos antes da forte contra o ventricular.\cite{guyton}.

O nodo atrioventricular tem como fun o retardar a passagem do impulso dos  trios para os ventriculos. No periodo que ocorre a ativa o do nodo atrioventricular, o impulso pecorre at  o  o feixe de His-Purkinje que os divide em dois ramos direito e esquerdo de His e se ramificam pela as fibras de Purkine, respons vel pela despolariza o e contra o do mioc rdio ventricular.  

A Figura \ref{fig:card} apresenta a caracter stica dos eventos fornecidos pelo os ciclos card acos, para o lado esquerdo do cora o. No ciclo card aco   poss vel extrair informa o de press o a rtica, press o atrial, press o ventricular, volume ventricular, fonocardiograma e por fim o eletrocardiograma que ser  a base desse trabalho.

\begin{figure}[H]
\begin{center}
			\caption{Eventos do Ciclo Card aco}
			\includegraphics[width=.5\textwidth]{Figuras/ecc.png}
            \vspace*{\fill} 
            \begin{quote} 
            \centering 
            Fonte: \cite{guyton}
            \end{quote}
            \vspace*{\fill}
			\label{fig:card}
\end{center}
\end{figure}


\subsection{O ELETROCARDIOGRAMA}

\hspace*{0.8cm}O eletrocardiograma   um sinal obtido atrav s da aquisi o dos sinais el tricos localizados no musculo card aco. A aquisi o do sinal   obtida atrav s da distribui o de sensores em ponto especifico sobre a pele utilizando eletrodos. O registro da atividade card aca   fornecido atrav s do eletrocardi grafo que fornece e registra gr ficos dos impulsos el tricos que estimulam a contra o do cora o em fun o do tempo \cite{thaler}.

A Figura \ref{fig:ecg} mostra um eletrocardiograma t pico com a identifica o das diversas ondas, tradicionalmente denominadas P, Q, R, S e T. A onda P   a primeira onda do sinal do eletrocardiograma e representa a despolariza o dos  trios direito e esquerdo sua amplitude   normalmente menor que 300$\mu$V, e a sua dura o   inferior a 120ms. Essa onda possui espectro abaixo de 10-15 Hz \cite{sornmo}.

\begin{figure}[!htb]
\begin{center}
			\caption{Caracter stica do sinal do eletrocardiograma}

\includegraphics[width=.5\textwidth]{Figuras/ecg1.PNG}
            \vspace*{\fill} 
            \begin{quote} 
            \centering 
            Fonte: \cite{sornmo}
            \end{quote}
            \vspace*{\fill}
			\label{fig:ecg}
\end{center}
\end{figure}



O complexo QRS   caracterizado pela despolariza o dos ventr culos esquerdo e direito em torno de 70ms a 110ms. Essa composi o de sinal   analisada atrav s das ondas Q, R e S. Esse complexo possui a maior amplitude em torno de 2mV a 3mV, e   a forma de onda que   analisada primeiro em aspectos computacionais devido a sua amplitude ser maior que as outras  \cite{sornmo}. 

O segmento ST   caracterizado pela a despolariza o dos ventr culos direito e esquerdo. Esse segmento inicia no final da onda S conhecido como ponto J at  a extremidade da onda T. J  o intervalo PQ corresponde ao tempo necess rio para que o impulso el trico se propague a partir do n dulo sino-atrial para os ventr culos   dependente da frequ ncia card aca \cite{sornmo}.

O intervalo RR corresponde ao comprimento de um ciclo card aco que representa o intervalo entre dois batimentos card acos e   usado para determinar arritmias irregulares. O intervalo QT apresenta o tempo do come o da despolariza o at  a repolariza o dos ventr culos. Este intervalo varia com a frequ ncia card aca o qual em taxas mais r pidas o mesmo torna-se reduzido. Por fim, a onda T representa a repolariza o ventricular com dura o de 300ms e   iniciada ap s o complexo QRS e as vezes   seguida por onda conhecida como U. A caracter stica   arredonda e est  associada a um  nico pico positivo \cite{sornmo}.


\subsection{DERIVA ES ELETROCARDIOGR FICAS}
\hspace*{0.8cm} A deriva o   uma linha que une eletricamente os eletrodos ao aparelho de eletrocardiograma. As deriva es eletrocardiogr ficas s o divididas em deriva o bipolar, subdivididas em deriva o I, deriva o II e deriva o III e as deriva es unipolares que s o divididas em V1, V2, V3, V4, V5, V6 as como aVR, aVL, aVF.

\subsubsection{As Tr s Deriva es Bipolares}

\hspace*{0.8cm}A caracteriza o do vetor resultante da ativa o el trica card aca   feita atrav s de par metros utilizando pares de eletrodos situados em pontos diferentes do corpo humano. Cada par metro resultante das medi es   chamado de deriva o, representando assim uma proje o desse vetor e espelhando a atividade el trica card aca em uma regi o do cora o \cite{webster}. Na aquisi o de sinais podem ser usadas as deriva es unipolar e bipolar.

As deriva es bipolares do plano frontal (DI, DII e DII) s o formadas pela a combina o de um par de eletrodos, sendo um posicionado nos membros inferiores e superiores depedendo da deriva o conforme a Tabela 1.


\begin{table}[H]
\centering
\caption{Posicionamento dos terminais}
\begin{center}
\begin{tabular}{|l|l|l|}
\hline
\textbf{Deriva es Bipolares} & \textbf{Terminal +} & \textbf{Terminal -} \\ \hline
\textbf{Deriva o I}          & Bra o esquerdo      & Bra o direito       \\ \hline
\textbf{Deriva o II}         & Perna esquerda      & Bra o direito       \\ \hline
\textbf{Deriva o III}        & Perna esquerda      & Bra o esquerdo      \\ \hline
\end{tabular}
\end{center}
 \vspace*{\fill} 
            \begin{quote}
            \centering 
            \vspace{12pt}Fonte: Elaborada pela autora 
            \end{quote}
            \vspace*{\fill}
			\label{fig:ramcor}
\end{table}

Na Figura \ref{fig:d1} temos o posicionamento espec fico de cada terminal (positivo e negativo) do eletrocardi grafo. Para a primeira deriva o o terminal negativo   conectado no bra o direito, enquanto o negativo posicionado no bra o direito. A deriva o II   resultante da combina o do terminal positivo alojado na perna esquerda e o terminal negativo alojado no bra o direito e a  ltima deriva o   consequente ao alojamento do terminal positivo na perna esquerda enquanto o negativo no bra o esquerdo. As deriva es bipolares de cada resultante da deriva o podem ser observadas na  Figura \ref{fig:d2}. O primeiro gr fico representa a deriva o I, o segundo a deriva o II e por fim a deriva o III.

\begin{figure}[!htb]
\begin{center}
			\caption{Deriva es bipolares no plano frontal}
			\includegraphics[width=.5\textwidth]{Figuras/d1.PNG}
            \vspace*{\fill} 
            \begin{quote} 
            \centering 
            Fonte: \cite{ebah}
            \end{quote}
            \vspace*{\fill}
			\label{fig:d1}
\end{center}
\end{figure}




\begin{figure}[H]
\begin{center}
			\caption{Comportamento da forma de onda da Deriva o I, II e III}
			\includegraphics[width=.5\textwidth]{Figuras/d2.PNG}
            \vspace*{\fill} 
            \begin{quote} 
            \centering 
            Fonte: \cite{guyton}
            \end{quote}
            \vspace*{\fill}
			\label{fig:d2}
\end{center}
\end{figure}
\subsubsection{Deriva es Tor cicas}
\hspace*{0.8cm} As deriva es tor cicas conhecida tamb m como deriva es precordiais s o obtidas atrav s do posicionamento de seis eletrodos situados na parede anterior ao t rax diretamente no cora o. O posicionamento dos eletrodos geram as seis deriva es que s o nominadas como deriva o V1, V2, V3, V4, V5 e V6, as respectivas deriva es medem a diferen a de potencial (ddp) entre o torax e o centro el trico do cora o. 
As deriva es v o desde de V1 localizado no espa o intercostal direito at  a deriva o V6, localizado no espa o intercostal esquerdo sendo mostrado na Figura \ref{fig:dop}.

\begin{figure}[H]
\begin{center}
			\caption{Posicionamentos dos seis eletrodos na regi o tor cica}
			\includegraphics[width=1\textwidth]{Figuras/v6.PNG}
            \vspace*{\fill} 
            \begin{quote} 
            \centering 
            Fonte: \cite{ronald}
            \end{quote}
            \vspace*{\fill}
			\label{fig:dop}
\end{center}
\end{figure}


Na Figura \ref{fig:d2}   poss vel observar a forma de onda de cada deriva o do plano horizontal na ordem de V1 a V6. O eletrodo da deriva o V1 avalia o lado direito do cora o, localizado no quarto espa o intercostal direito, justaestemal, a deriva o V2 avalia o lado direito do cora o localizado no quarto espa o intercostal esquerdo, justamental, a deriva o V3 avalia uma regi o intermedi ria entre V2 e V4, a deriva o V4 localiza-se no quinto espa o intercostal esquerdo na linha hemiclavicular, a deriva o V5 entre o 5  espa o intercostal esquerdo na linha axilar anterior e a deriva o V6 localiza-se no quinto espa o intercostal esquerdo na linha axilar m dia \cite{ronald}.

\begin{figure}[H]
\begin{center}
			\caption{Comportamento das forma de ondas da Deriva o V1, V2, V3, V4, V5 e V6.}
			\includegraphics[width=.5\textwidth]{Figuras/DEV.PNG}
            \vspace*{\fill} 
            \begin{quote} 
            \centering 
            Fonte: \cite{guyton}
            \end{quote}
            \vspace*{\fill}
			\label{fig:d2}
\end{center}
\end{figure}


\subsubsection{As Deriva es Unipolares Aumentadas}


\hspace*{0.8cm}As deriva es  unipolares obt m potenciais locais, ou seja, cada deriva o   obtida atrav s do posicionamento de eletrodos exploradores. A deriva o VR   obtida atrav s do alojamento do eletrodo explorador no bra o direito, a deriva o VL   adquirida posicionado o eletrodo no bra o esquerdo e quando o eletrodo   colocado na perna esquerda tem-se a deriva o VF que s o baseadas em um ponto conhecido como Terminal Central de Wilson (WCT), conforme apresentado na Figura \ref{fig:ab}.

As deriva es unipolares aumentadas denominadas como aVR, aVL e AVF s o fornecidas atrav s de dois membros s o conectados ao terminal negativo atrav s de resist ncias el tricas e o terceiro membro conectado ao terminal positivo.

\begin{figure}[H]
\begin{center}
			\caption{Terminal Central de Wilson}
			\includegraphics[width=.5\textwidth]{Figuras/WCT.PNG}
            \vspace*{\fill} 
            \begin{quote} 
            \centering 
            Fonte: \cite{webster}
            \end{quote}
            \vspace*{\fill}
			\label{fig:ab}
\end{center}
\end{figure}
O tipo de configura o de resist ncia mostrada na Figura \ref{fig:ab} tende a reduzir a amplitude do sinal. Para o aumento da amplitude o  circuito   modificado para a gera o da deriva o unipolar aumentada (aVR, aVL e aVF) que s o as deriva es que s o mostradas em monitores card acos.

A Figura \ref{fig:aca} representa a configura o do terminal de Wilson com a resist ncia el trica modificada para gerar o aumento da amplitude. Neste tipo de configura o, dois membros s o conectados, por meios das resist ncias el tricas, ao terminal negativo e o terceiro membro no terminal positivo para a deriva o aVL o terminal positivo precisa est  conectado ao bra o esquerdo e a ultima deriva o aumentada denominada AVF precisa que a perna esquerda esteja conectado ao terminal positivo \cite{guyton}. 

\begin{figure}[!htpb]
\begin{center}
			\caption{Deriva es unipolares a) aVR, b)aVL, c) aVF }
			\includegraphics[width=.5\textwidth]{Figuras/avr.PNG}
            \vspace*{\fill} 
            \begin{quote} 
            \centering 
            Fonte: \cite{webster}
            \end{quote}
            \vspace*{\fill}
			\label{fig:aca}
\end{center}
\end{figure}

Para obter a deriva o aVR, a resist ncia conectada ao eletrodo do bra o direito   modificada enquanto os demais continuam com o mesmo valor. Para a gera o da deriva o aVL a resist ncia que   modificada fica conectada ao eletrodo do bra o esquerdo e para a deriva o aVF a resist ncia   modificada no eletrodo que   conectado a perna esquerda. Esse arranjo serve para aumentar a amplitude de cada deriva o unipolar aumentada (aVR, aVL e aVF), essas deriva es tem a caracter stica de formas de ondas mostrada na Figura \ref{fig:aba}, a primeira a forma de onda representa a primeira unipolar aumentada aVR, a segunda aVL e a terceira aVF.

\begin{figure}[!htpb]
\begin{center}
			\caption{Registro das tr s deriva es unipolares aumentadas}
			\includegraphics[width=.5\textwidth]{Figuras/tres.PNG}
            \vspace*{\fill} 
            \begin{quote} 
            \centering 
            Fonte: \cite{guyton}
            \end{quote}
            \vspace*{\fill}
			\label{fig:aba}
\end{center}
\end{figure}

Na Tabela 2 apresenta as rela es matem ticas para obten o das deriva es eletrocardiogr ficas a partir do posicionamento dos eletrodos  posicionandos nos membros inferiores e superiores (LA - bra o esquerdo, LL - bra o direito, RA- perna esquerda e RL- perna direita) e dos eletrodos posicionando na regi o tor cica (v1, v2, v3, v4, v5 e v6).  

\begin{table}[H]
\centering
\caption{Representa o Algebrica atrav s do posicionamento dos eletrodos }
\begin{center}
\begin{tabular}{|l|l|l|ll}
\cline{1-3}
Tipo de Deriva o                                                                       & Eletrodos Utilizados   & Defini o                                                                                                                                                             &  &  \\ \cline{1-3}
\begin{tabular}[c]{@{}l@{}}Bipolar \\ ou \\ Deriva o de membros Einthoven\end{tabular} & LA, RA, LL, RL         & \begin{tabular}[c]{@{}l@{}}I= LA-RA\\ II = LL-RA\\ III = LL-LA\end{tabular}                                                                                           &  &  \\ \cline{1-3}
Aumentadas                                                                              & LA, RA, LL, RL         & \begin{tabular}[c]{@{}l@{}}aVR= RA-1/2(LA+LL)\\ aVL =LA-1/2(LL+RA)\\ aVF = LL-1/2(LA+RA)\end{tabular}                                                                 &  &  \\ \cline{1-3}
Unipolar                                                                                & V1, V2, V3, V4, V5, V6 & \begin{tabular}[c]{@{}l@{}}V1 = v1-(RA+LA+LL)/3\\ V2=v2-(RA+LA+LL)/3\\ V3=v3-(RA+LA+LL)/3\\ V4=v4-(RA+LA+LL)/3\\ V5=v5-(RA+LA+LL)/3\\ V4=v6-(RA+LA+LL)/3\end{tabular} &  &  \\ \cline{1-3}
\end{tabular}
\end{center}

 \vspace*{\fill} 
            \begin{quote}
            \centering 
            \vspace{12pt}Fonte: Elaborada pela autora 
            \end{quote}
            \vspace*{\fill}
			\label{fig:ramcor}
\end{table}            

\subsection{O TRI NGULO DE EINTHOVEN }

\hspace*{0.8cm}O tri ngulo de Einthoven   uma representa o vetorial em torno do cora o que permite estudar as for as eletromotrizes geradas pelo cora o. Esse representa o fornece  diferen a de potencial entre eletrodos localizados em diversos membros sobre a pele conhecidas como deriva o I, II e III, caso os potenciais el tricos de duas das tr s deriva es bipolares forem conhecidos e a terceira pode ser determinada matem ticamente a partir das duas primeiras. Na Figura \ref{fig:jup} temos o tri ngulo de Einthoven.

\begin{figure}[!htpb]
\begin{center}
			\caption{Tri ngulo de Einthoven.}
			\includegraphics[width=.5\textwidth]{Figuras/triangulo.png}
            \vspace*{\fill} 
            \begin{quote} 
            \centering 
            Fonte: \cite{souza}
            \end{quote}
            \vspace*{\fill}
			\label{fig:jup}
\end{center}
\end{figure}

A deriva o I registra a diferen a de potencial entre o bra o direito e bra o esquerdo, a deriva o II entre o bra o direito e perna esquerda e deriva o III entre o bra o esquerdo e a perna direita. 

A Figura \ref{fig:jup} representa a figura geom trica do Tri ngulo de Einthoven. O braco direito e esquerdo e a perna esquerda formam os  pices do tri ngulo que envolve o cora o. Os  pices superiores representam os pontos localizados no bra o esquerdo e bra o direio e o  pice inferior representa a perna esquerda se conecta aos dois pontos, a liga o entre todos os apices fornecem o Tri ngulo de Einthoven.

\subsection{ELETRODOS}

\hspace*{0.8cm}Os eletrodos s o sensores que t m como fun o medir e registrar os biopotenciais atrav s da interface fornecida pelo corpo humano e o circuito de aquisi o. Os eletrodos podem ser classificados como eletrodos invasivos e n o invasivos. Os eletrodos invasivos s o inseridos internamente ao corpo para a detec o de biopotenciais, e tamb m s o conhecidos como subcut neos. Os eletrodos n o invasivos s o utilizados para detectar biopotenciais a partir da superf cie do corpo \cite{vier}.

Normalmente, os eletrodos mais utilizados para monitora o do eletrocardiograma s o do tipo autoadesivos constitu dos por gel s lido de cloreto de pot ssio, uma capa pl stica para manter a umidade do gel, pino de a o inox, contra pino de prata cloreto de pot ssio, papel protetor sendo o contato el trico com a pele \cite{Site22}. Na Figura \ref{fig:eletrodo} temos um exemplo de eletrodo do tipo n o invasivo, comum em sistemas de ECG.
%http://www.fibracirurgica.com.br/eletrodo-descartavel-espuma-2223brq-pacote-c-50-unidades-3m/p

\begin{figure}[!htpb]
\begin{center}
			\caption{Eletrodos do tipo n o invasivo.}
			\includegraphics[width=.5\textwidth]{Figuras/eletrodo.png}
            \vspace*{\fill} 
            \begin{quote} 
            \centering 
            Fonte: Elaborada pela autora
            \end{quote}
            \vspace*{\fill}
			\label{fig:eletrodo}
\end{center}
\end{figure}

\subsection{INTERFER NCIA EM SINAIS DE ECG}

\hspace*{0.8cm}O sinal ECG   sens vel a interfer ncias, principalmente devido a caracter stica deste sinal possui pequena amplitude, por volta de 1mVpp. No sistema de capta o do sinal ECG utiliza no m nimo dois eletrodos explorat rios e um de refer ncia.

O tipo de interfer ncia mais comum em sistema de aquisi o de sinais biol gicos   oriundo a frequ ncia da rede el trica e induzido por meios de acoplamentos entre o comum do circuito eletr nico, ou terra, e a rede el trica que alimenta o dispositivo  pode variar dependendo do pa s sendo  50Hz e 60Hz.

A interfer ncia pode ser inserida ao sinal devido a diferen a de imped ncia entre os eletrodos, o que resulta em desbalanceamento nas tens es de modo comum e as converte em tens es diferenciadas, sendo assim amplificadas. Esse fen meno   nomeado de efeito de divisor de potencial.\cite{webster}

A Figura \ref{fig:tri}  representa a interfer ncia de 60 Hz no sinal do eletrocardiograma. O sinal em azul representa em condi es ideais e o vermelho com a interfer ncia da frequ ncia.

\begin{figure}[!htb]
\begin{center}
			\caption{Sinal ECG com interfer ncia devido a rede el trica}
			\includegraphics[width=.5\textwidth]{Figuras/freq.PNG}
            \vspace*{\fill} 
            \begin{quote} 
            \centering 
            Fonte: \cite{android}
            \end{quote}
            \vspace*{\fill}
			\label{fig:tri}
\end{center}
\end{figure}

Para solucionar problemas com interfer ncias no sinal do ECG s o aplicados v rios tipos de filtros sendo filtro passa- baixa, passa -- alta, notch e casado.
\begin{itemize}
\item Filtro passa- alta: Aplicado para remo o do sinal base.
\item Filtro passa -- baixa: Remo o do ru do muscular.
\item Filtro Notch:  Remo o do ru do de 60Hz da rede el trica.
\item Filtro Casado: Utilizado para comparar uma forma de onda com outra utilizada como onda padr o. Esse tipo de filtro   fundamental para que seja encontrado o completo QRS e seja, portanto, uma an lise para poss vel detec o de anomalia. 

\end{itemize}


%Chamda de site \cite{wiki}.

\subsection{CIRCUITO DE AQUISI O}
\hspace*{0.8cm} O circuito de aquisi o   respons vel em adquirir o sinal de ECG. A aquisi o do sinal   de extrema import ncia, pois   a primeira etapa do sistema de capta o e transmiss o de sinais de ECG. Este circuito   composto basicamente por um amplificador de instrumenta o espec fico para sinais de ECG, devido a sua baixa energia \cite{pds}.


O amplificador de instrumenta o   um circuito integrado e geralmente constitu do por tr s amplificadores operacionais, sendo aplicado para medir sinais anal gicos de baixa amplitude que requer exatid o. A Figura \ref{fig:cir} apresenta a estrutura deste amplificador. Pode se observar que este possui dois est gios: um est gio de entrada, constitu do por dois amplificadores n o-inversores e o segundo est gio formado por um amplificador subtrator \cite{veloso}. 

\begin{figure}[!htb]
\begin{center}
			\caption{Estrutura caracter stica de um amplificador de instrumenta o.}
			\includegraphics[width=.5\textwidth]{Figuras/circuito.png}
            \vspace*{\fill} 
            \begin{quote} 
            \centering 
            Fonte: \cite{sugi}
            \end{quote}
            \vspace*{\fill}
			\label{fig:cir}
\end{center}
\end{figure}

Ap s a aquisi o do sinal, o sinal de ECG pode ser enviado para outros tipos de circuitos a depender do sistema, como por exemplo, circuitos de amplifica o de segundo est gio e filtragem e circuitos de modula o.

\subsection{FILTROS ATIVOS}
\hspace*{0.8cm}Os filtros s o circuitos eletr nicos que executam o processamento de sinais especificamente atenuando frequ ncias indesejadas de sinais de entradas de um circuito, tem como caracter stica permitir a passagem ou suprimir sinais com certas frequ ncias dependendo da topologia empregada. Um circuito de filtro ativo pode ser constru do utilizando componentes tais como: amplificadores operacionais, transistores, resistores, capacitores.

Um filtro que permite a passagem de sinais de tens es e corrente somente em frequ ncias abaixo de certo limite, atenuando os sinais cuja a frequ ncia ultrapassar esse valor   chamado de filtro passa-baixa. Um filtro que permite a passagem de sinais acima de uma frequ ncia de corte   um filtro passa-alta. Quando um circuito apresenta a passagem de sinal acima de uma frequ ncia de corte e abaixo de outra frequ ncia de corte,   chamado de filtro passa-banda \cite{boy}.

\subsection{CONVERSOR ANAL GICO/ DIGITAL}

\hspace*{0.8cm}Um conversor anal gico-digital recebe uma tens o anal gica de entrada e posteriormente produz um sinal digital de sa da correspondente a entrada anal gica. Para convers o do sinal anal gico   realizado a amostragem do sinal. As principais t cnicas de convers o s o a paralelo (tamb m conhecida como flash), contador de rampa cont nuo e aproxima o sucessivas. Suas caracter sticas desejadas est o na precis o e tempo m ximo de convers o \cite{digi}.

Alguns circuitos integrados (CIs) implementam conversores anal gico-digital e podem ser utilizados utilizando algum tipo de protocolo. Um deles   o MCP3008, que   um CI de 16 pinos, sendo 8 entradas anal gicas com 6 entradas de sinais de controle, com sa da utilizando o protoloco  de comunica o SPI (\textit{Serial Peripheral Interface}). Esse protocolo permite a comuniza o utilizando a interconex o de dispositivos ou chips de forma que estes troquem informa o entre si.
%http://www.alan.eng.br/disc_microprocessadores/pratica8_spi.pdf


\subsection{CLIENTE-SERVIDOR}

\hspace*{0.8cm}Os dipositivos em geral, utilizam protocolos de comunica o, sendo os principais: IP (\textit {Internet Protocol}), TCP (\textit {Transmission Control Protocol}) e HTTP, (\textit{Hypertext Transfer Protocol}). Estes s o protocolos de padroniza o que possibilitam a troca de informa o entre si de diferentes equipamentos. Esses protocolos s o usados na transfer ncia de informa es e s o tamb m aplicado na arquitertura cliente-servidor.

\hspace*{0.8cm}Esse tipo de arquitetura   caracterizada em dois n vel chamada cliente-servidor, e funciona da seguinte maneira: o cliente socilita um recurso e o servidor responde com a solicita o do cliente, utilizando  a internet como principal meio de comunica o. A Figura \ref{fig:cli} representa a estrutura do processo cliente servidor, no qual o cliente solicita e o servidor responde a solicita o do cliente.

O termo cliente-servidor refere-se quando se tem dois ou mais computadores interagindo de modo que um oferece servi o a outro, permite ao usu rio acessar a informa o e servi os em qualquer lugar. 
\begin{itemize}
\item Cliente: Conhecido como ``front-end'' ou ``WorkStation''   um recurso o qual usu rio interage com o usu rio atrav s de uma interface gr fica, ou n o, permitindo consultas de dados, manipula o  de tela, entrada e valida o de dados, gerenciamento de som e  udio e entre outras fun es. Normalmente o cliente   dedicado a sess o do usu rio solicitando servi os, iniciando e finalizando a sess o.
\item Servidor: Respondem as solicita es gerada pela o cliente,   um processo cont nuo, oferecendo servi os quando solicitado pelo o cliente,integra servi os de armazenamento, acesso, organiza os dados, atualiza os dados, gerencia os recursos compartilhados e entre outros servi os. 

\end{itemize}

\begin{figure}[H]
\begin{center}
			\caption{Estrutura Cliente-Servidor}
			\includegraphics[width=.9\textwidth]{Figuras/servidor.PNG}
            \vspace*{\fill} 
            \begin{quote} 
            \centering 
            Fonte: \cite{tanen}
            \end{quote}
            \vspace*{\fill}
			\label{fig:cli}
\end{center}
\end{figure}


\newpage

\section{M DULO DE AQUISI O E TRANSMISS O DE SINAIS ELETROCARDIOGR FICO}

\hspace*{0.8cm}Este cap tulo aborda o desenvolvimento do m dulo de aquisi o de sinais ECG, como tamb m o circuito de convers o anal gico e digital e tranmiss o desses dados via Wi-Fi para um computador. Al m disso, esse cap tulo aborda aspectos da placa de aquisi o de sinais, desenvolvimento do software para transmiss o do sinal, respons vel por receber os dados e exibi-los na tela do computador. 

Uma vez que o sinal de ECG possui a amplitude em escala de milivolts   necess rio um mecanismo para que seja realizado o condicionamento deste sinal. O circuito respons vel por tratar o sinal   composto de amplificadores, filtros e conversores A/D acoplados com o Raspberry Pi 3. 

O desenvolvimento do projeto consistiu em duas etapas: a primeira etapa do projeto de pesquisa, consistiu na aquisi o, leitura e  teste do hardware, composto pela a Placa End Front-End e o MCP3008. Nesta estapa foram parametrizados as limita es dos componentes, tratamento e visualiza o dos sinais que comp e o eletrocardiograma. Nesta fase utilizou-se para a simula o o simulador ECG 12 deriva es HS-15, da RQ Mediq, que representam a c pia aut ntica do sinal gerado pelo o corpo humano.


A segunda parte do projeto consistiu na transfer ncia dos dados para um computador, utilizando o protocolo de comunica o em rede (TCP/IP). Nesta etapa foram definidos os protocolos e a linguaguem de programa o, com a finalidade de transferir  as amostras geradas pelo o simulador de sinais ao computador, transpondo dessa maneira as amostras em um determinado espa o de tempo em forma de gr fico, tendo como caracter stica principal sinais digitais de ECG compostos pelas as ondas P, complexo QRS e a onda T. 

A Figura \ref{fig:pap} representa as etapas desde a aquisi o dos sinais atrav s dos eletrodos, cabos e circuito de aquisi o at  a transfer ncia ao usu rio utilizando conversores, linguagem de programa o embarcado no Raspberry Pi. 
    
\begin{figure}[H]
\begin{center}
			\caption{Diagrama do trabalho  desenvolvido.}
			\includegraphics[width=.9\textwidth]{Figuras/pac.PNG}
            \vspace*{\fill} 
            \begin{quote} 
            \centering 
            Fonte: Elaborada pela autora
            \end{quote}
            \vspace*{\fill}
			\label{fig:pap}
\end{center}
\end{figure}

Lista de materiais do projeto


\begin{itemize}
\item 01 AD8232
\item 01 Simulador de ECG Handy Sim HS-15
\item 01 MCP3008
\item 01 Raspberry Pi 3
\item 01 Cabo HDMI
\item 01 Cabo paciente de 03 vias.
\item 03 eletrodos auto adesivos
\item 01 Carregador 5,1 volts/2,5 Ampere
\item 01 Computador
\item 01 CABO PACIENTE ECG 5 VIAS AE-0H002-0
\item Tornado vers o: 4.5.2
\end{itemize}

\subsection{M DULO DE AQUISI O DE SINAIS ECG}
\hspace*{0.8cm} O m dulo de aquisi o de sinais ECG   composto basicamente por um componente AD8232, constitu do basicamente por amplificador  instrumenta o, um filtro de condicionamento do sinal e um aplificador de realimenta o ou um circuito de perna direita cuja fun o desse componente   obter e filtrar os sinais de eletrocardiograma integrado em  nico capsulamento, comercializado pelo o fabricante Analog Devices.

\subsubsection{Placa ECG Front-End}
\hspace*{0.8cm}A placa de ECG escolhida   aplicada para  tratamento e condicionamento de pequenos sinais que comp em os sinais  ECG, projetado para filtrar, extrair e amplificar pequenos sinais card acos,   da empresa SparkFun, baseada no CI AD8232 da Analog Devices,   composto por 20 pinos e internamente este pinos s o entrada e sa da de amplificadores de instrumenta o e operacional aplicados a este componente, sendo detalhado na Figura \ref{fig:core}.

O CI AD8232 possibilita ao usu rio implementar v rias fun es como o uso de filtros para eliminar artefatos de movimentos, ou seja, anormalidades que n o resultam da atividade el trica do cora o, amplificadores operacional permite filtro passa-baixa para remover o ru do adicional, filtro passa alta para eliminar sinais provenientes do movimento muscular e do eletrodo.

A Figura \ref{fig:circ} representa a caracter stica externa e interna do AD8232 e as respectivas pinagens composto por 20 pinos. Internamente possui um amplificador de instrumenta o (IA), um amplificador operacional (A1), um aplificador de unidade perna direita (A2) e um buffer da entrada de refer ncia (A3).

A configura o interna da placa da Sparkfun   projetada para realizar o monitoramento da forma de onda composto pelas as ondas P, complexo QRS e a onda T. Esse sinal possui baixa amplitude logo   suscept vel a ru dos, portanto, para obter o sinal sem distor o   acoplado  filtro passa-alta com dois p los em 0,5 Hz seguido de um filtro passa-baixa com dois-p los em 40 Hz. Um terceiro circuito de alimenta o   conduzido para uma rejei o de modo comum.

Especifica o:
\begin{itemize}
\item Tens o de Opera o 3.V.
\item Sa da anal gica
\item Peso 4g
\item 03 Pinos de Entrada (RA, LA e RL)
\item Temperatura de Funcionamente 40 a 55.%N o seicolocar o grau
\item Led indicador
\item Tamanho 4mm x 4mm

\end{itemize}

\begin{figure}[H]
\begin{center}
			\caption{Pinagem do AD8232}
			\includegraphics[width=.9\textwidth]{Figuras/adc.PNG}
            \vspace*{\fill} 
            \begin{quote} 
            \centering 
            Fonte: \cite{analogic}
            \end{quote}
            \vspace*{\fill}
			\label{fig:core}
\end{center}
\end{figure}

A placa de ECG Front-End utilizada   da empresa Sparkfun e possui internamente o AD8232, resistores e capacitores de precis o. Os pinos RA, LA e RL s o respectivamente do bra o direito, bra o esquerdo e perna direita. Externamente possui um led indicador que pulsa na frequ ncia do ciclo card ado al m de outras 6 sa das que s o de aterramento do circuito (GND), tens o de alimenta o (3,3V), sa da anal gica (OUTPUT), detectores de eletrodo (LO-/LO+) e desligamento (SDN), demonstrado na Figura \ref{fig:circ}.

\begin{figure}[H]
\begin{center}
			\caption{Placa ECG}
			\includegraphics[width=.5\textwidth]{Figuras/pc.PNG}
            \vspace*{\fill} 
            \begin{quote} 
            \centering 
           Fonte: \cite{analogic}
            \end{quote}
            \vspace*{\fill}
			\label{fig:circ}
\end{center}
\end{figure}
Esse componente n o possui conversor anal gico digital, por m o mesmo permite que integre outros componentes atrav s da sa da (OUTPUT). O componente escolhido e utilizado para converter o sinal em digital   chamado de MCP3008 que ser  abordado na Subse o \ref{sec:conversorAD}.

\subsubsection{Conversor Anal gico Digital}\label{sec:conversorAD}

\hspace*{0.8cm}O MCP3008   um conversor anal gico digital program vel que opera na faixa de 2,7V a 5,5V. A comunica o com os dispositivos   realizada usando
interface serial simples compat vel com o protocolo SPI.

Esse componente possui 8 entradas anal gicas do canal CH0 ao CH7 e 8 pinos configur veis. Esses pinos s o conhecidos como VDD, VREF, AGND, CLK, DOUT, DIN, CS/SHDN e DGND, detalhado na Tabela 1. Os pinos DIN, DOUT e CS/SHDN pode ser conectados a outro dispositivos com interface SPI para realizar a convers o de anal gico e digital. 

A Figura \ref{fig:cadc} representa a caracter stica externa do MCP3008 com as respectivas pinagens CH0 a CH7 e entradas anal gicas: VDD, VREF, AGND, CLK, DOUT, DIN, CS/SHDN E DGND.

\begin{table}[]
\centering
\begin{center}
\caption{Descri o dos Pinos}
\label{center}
\begin{tabular}{|l|l|l|}
\hline
\multicolumn{1}{|c|}{\textbf{PINOS}} & \multicolumn{1}{c|}{\textbf{SIMBOLOGIA}} & \multicolumn{1}{c|}{\textbf{DESCRI O}} \\ \hline
1                                    & CH0                                      & Entrada Anal gica                       \\ \hline
2                                    & CH1                                      & Entrada Anal gica                       \\ \hline
3                                    & CH2                                      & Entrada Anal gica                       \\ \hline
4                                    & CH3                                      & Entrada Anal gica                       \\ \hline
5                                    & CH4                                      & Entrada Anal gica                       \\ \hline
6                                    & CH5                                      & Entrada Anal gica                       \\ \hline
7                                    & CH6                                      & Entrada Anal gica                       \\ \hline
8                                    & CH7                                      & Entrada Anal gica                       \\ \hline
9                                    & DGND                                     & \textit{Digital Ground}                 \\ \hline
10                                   & CS/SHDN                                  & \textit{Chip Select/Shutdown Input}     \\ \hline
11                                   & DIN                                      & \textit{Serial Data In}                 \\ \hline
12                                   & DOUT                                     & \textit{Serial Data Out}                \\ \hline
13                                   & CLK                                      & \textit{Serial CLK}                     \\ \hline
14                                   & AGND                                     & \textit{Analog Ground}                  \\ \hline
15                                   & VREF                                     & Tens o de Refer ncia                    \\ \hline
16                                   & VDD                                      & + 2,5V a 5,5V                           \\ \hline
\end{tabular}
\end{center}
 \vspace*{\fill} 
            \begin{quote}
            \centering 
            \vspace{12pt}Fonte: Elaborada pela autora 
            \end{quote}
            \vspace*{\fill}
			\label{fig:ramcor}
\end{table}


\begin{figure}[H]
\begin{center}
			\caption{Conversor ADC}
			\includegraphics[width=.5\textwidth]{Figuras/mcp3008.PNG}
            \vspace*{\fill} 
            \begin{quote} 
            \centering 
           Fonte: \cite{analogica}
            \end{quote}
            \vspace*{\fill}
			\label{fig:cadc}
\end{center}
\end{figure}

Caracter stica:

\begin{itemize}
\item ADC de 10 bits
\item 8 canais
\item Interface SPI
\item Taxa de amostragem: 200kSPS
\item Tens o de alimenta o 2,7 V a 5,5 V
\item DIP de 16 pinos
\item Tempo de convers o %10$\mu$s
\item Corrente de suprimento: %425$\mu$A

\end{itemize}

\subsubsection{Raspberry Pi 3}

\hspace*{0.8cm}O Raspberry Pi   um microcomputador com tamanho equivalente a um cart o de cr dito, com hardware livre, criado com o intuito de disseminar o ensino de programa o para jovens e crian as. Seu custo est  em torno de 10 a 35 d lares e   desenvolvido pela Raspberry Pi Foundation, uma institui o sem fins lucrativos do Reino Unido.
% https://www.raspberrypi.org 

Em compara o a outros sistemas similares como o Arduino, suas vantagens est o em oferecer ao usu rio a possibilidade de ser programada em C++, Java e Phyton, enquanto o Arduino   baseado em C++.  Al m disso, permite ao usu rio executar v rios programas ao mesmo tempo enquanto no Arduino s    poss vel apenas um por vez. O modelo que ser  utilizado neste trabalho   o  ltimo lan amento conhecido como Raspberry Pi 3 Model B, devido ao adaptador wireless integrado, o que permite ao usu rio conex o   rede sem fio, possibilitando diversos tipos de aplica es \cite{usp}. 

Esse microcomputador   um computador em qualquer uma de suas vers es existente no mercado, portanto   necess rio instalar um Sistema Operacional (SO), como Linux ou Windows e recomenda-se que seja instalado a vers o Debian Linux, por m existe outras vers es dispon veis.  O sistema operacional   armazenado no cart o de mem ria SD, o qual funciona como disco r gido, localizado na parte inferior do hardware \cite{usp}.

A Figura \ref{fig:raspb} representa a vista superior da Raspberry Pi 3 e seus perif ricos que permite gerenciar teclado, monitor, mouse, pen drive e entre outros perif ricos utilizando as 4 portas USB. Possui 40 conectores de GPIO, conector Ethernet, sa da de  udio e v deo, conector de c mera, sa da de v deo, comunica o com redes e formas de comunica o serial e entre outros.

\begin{figure}[H]
\begin{center}
			\caption{Estrututa externa do Raspberry Pi 3}
			\includegraphics[width=.9\textwidth]{Figuras/RASP.PNG}
		
            \vspace*{\fill} 
            \begin{quote} 
            \centering 
           Fonte: \cite{raspb}
            \end{quote}
            \vspace*{\fill}
			\label{fig:raspb}
\end{center}
\end{figure}


Caracter stica:

\begin{itemize}
\item Processador Bradcom BCM2837 64bit ARMv8 Cortex -- A53 Quad-Core.
\item Mem ria RAM: 1GB
\item Adaptador Wifi 802.11n Integrado
\item Bluetooth 4.1 BLE integrado
\item Conector de video HDMI
\item 4 Portas USB 2.0
\item Conector Ethernet 
\item Interface para c mera (CSI)
\item Interface para display (DSI)
\item Slot para cart o SD
\item Conector para  udio e v deo
\item GPIO de 40 pinos
\item Dimens o: 80 x 56 x 17mm
\end{itemize}

O Raspberry Pi n o possui identifica o dos pinos na placa, logo   necess rio acessar o mapeamento dos pinos de acordo com Figura \ref{fig:pinosraspb}. Os pinos 2 e 4 representam pinos de alimenta o de 5 Volts e especificamente essas portas n o podem entrar em contato com as demais diferentemente do pino 1 e 17 que possuem sa da 3,3V,  possui tamb m 8 portas (8, 25, 6, 14, 20, 30 e 34) de aterramento (GROUND). Os pinos 3 e 4 s o destinado ao uso protocolo de comunica o I2C, sendo um destinado a dados serial e outra para CLK, os pinos 8 e 10 s o portas serial e utilizam protocolo RS-232 para transmiss o e recep o de sinal digital. Os pinos 19, 21 ,23, 24, 26 permitem que a raspberry se comunique de forma serial Full Duplex s ncrono com outro dispositivo perif rico de forma bidirecional, por m o protocolo SPI precisa est  implementado. As portas 26 e 28 s o  portas do ID EEPROM  (\textit{Electrically-Erasable Programmable Read-Only Memory}), ou seja   um tipo de mem ria que pode ser programada e apagada v rias vezes, atrav s de uma tens o el trica interna ou externa e as demais portas s o pinos de GPIO que fazem envio e transmiss o de sinais digitais. %ref [http://blog.fazedores.com/raspberry-pi-b-introducao-porta-gpio/]

\begin{figure}[H]
\begin{center}
			\caption{Mapeamento dos pinos do Raspberry Pi 3}
			\includegraphics[width=.9\textwidth]{Figuras/pins.PNG}
			\vspace*{\fill} 
            \begin{quote} 
            \centering 
           Fonte: \cite{ras}
            \end{quote}
            \vspace*{\fill}
			\label{fig:pinosraspb}
\end{center}
\end{figure}

%fonte: https://jeremylindsayni.wordpress.com/2017/04/05/turning-gpio-pins-high-and-low-on-a-raspberry-pi-3-using-net-core-2-and-ubuntu/
\subsubsection{Cabo Paciente}
\hspace*{0.8cm} O cabo paciente   utilizado em monitores cardiacos e eletrocardiografo que auxiliam esses equipamentos a monitorar sinais ECG. O cabo utilizado   fabricado pela a empresa Dixtal Philips o qual possui certifica o pela a ANVISA. O modelo utilizado para o experimento   definido como cabo paciente ECG 5 vias (AE-0H0020). Este modelo est  apresentado na Figura \ref{fig:cabopaciente}.
\begin{figure}[H]
\begin{center}
			\caption{Cabo paciente (AE-0H002-0) com as vias modificadas}
			\includegraphics[width=.5\textwidth]{Figuras/Cabopaciente.PNG}
 			\vspace*{\fill} 
            \begin{quote} 
            \centering 
           Fonte: Elaborada pela autora
            \end{quote}
            \vspace*{\fill}
			\label{fig:cabopaciente}
\end{center}
\end{figure}

\subsubsection{Simulador de Sinais Vitais}
\hspace*{0.8cm} No mercado existem diversos modelos de simulador de sinais vitais que simulam ondas no formato anal gico. A diferen a   que alguns possuem 5 pinos, ou seja   limitado para algumas deriva es e outros possuem 10 pinos, que geram todas as deriva es. O mais avan ado neste ramo   conhecido como PROSIM4- FLUKE, desenvolvido para gerar as 12 formas de deriva es ECG e oferece simula es com os par metros de marca-passo, teste de arritmia, respira o, press o invasiva, temperatura, d bito card aco, simula o de ECG fetal/materno. Por possuir diversas funcionalidades, seu pre o   elevado quando comparado a outros equipamentos. Este modelo   demonstrado na Figura \ref{fig:simvitais}.

\begin{figure}[H]
\begin{center}
			\caption{Simulador de sinais vitais - PROSIM3-FLUKE}
			\includegraphics[width=.4\textwidth]{Figuras/fluke.png}
			\vspace*{\fill} 
            \begin{quote} 
            \centering 
           Fonte: \cite{fluke}
            \end{quote}
            \vspace*{\fill}
			\label{fig:simvitais}
\end{center}
\end{figure}

O HS-15   o modelo utilizado nesse projeto para teste e gerar os sinais vitais em modo ECG. Este simulador permite alterar frequ ncia, amplitude e opera  em cinco modos que fornecem a forma de onda quadrada, triangular, seno, pmp e ECG.  O pulso de ECG pode ser modificado em 30-60-80-120-180-240 btm. 
	

A Figura \ref{fig:hs15} representa o simulador de sinais HS-15 da empresa Handy Sim, desenvolvido para ser utilizado em procedimento de verifica o de par metro, manunten o, avalia o de desempenho, teste e calibra o em monitores card acos e eletrocardi grafo. 
 

%http://catalogohospitalar.com.br/hs15-simulador-de-ecg-c-12-variacoes-para-eletrocardiografo-multicanal.html

\begin{figure}[H]
\begin{center}
			\caption{Simulador HS-15}
			\includegraphics[width=.4\textwidth]{Figuras/HS150.PNG}
			\vspace*{\fill} 
            \begin{quote} 
            \centering 
           Fonte: \cite{HS15}
            \end{quote}
            \vspace*{\fill}
			\label{fig:hs15}
  
\end{center}
\end{figure}

\newpage


\section{IMPLEMENTA O DO PROJETO}

\hspace*{0.8cm} Este cap tulo aborda o desenvolvimento do m dulo de aquisi o do sinal ECG incluindo o circuito de aquisi o de sinal, contemplando a convers o anal gico-digital e a transmiss o do sinal ECG. O diagrama esquem tico deste sistema est  exposto na  Figura \ref{fig:dia}. Deste diagrama, tem-se os eletrodos, que est o conectados no paciente; o cir\-cuito de aquisi o para tratar e acomodar o sinal enquanto o MCP3008   utilizando como conversor anal gico/digital e se conecta com o Raspberry Pi 3 atrav s da porta SPI. Ambos est o configurados para transmitir o sinal ECG. O sistema   iniciado com o pedido realizado pelo usu rio, utilizando a arquitetura cliente-servidor presente nos computadores. Internamente o cliente
faz a requisi o ao servidor presente no Raspberry Pi 3 e o interpreta, utilizando a rede Wi-Fi como meio de comunica o, realizando a leitura de dados da porta do MCP3008, a qual est  conectada com o circuito de aquisi o e por fim, o cliente recebe as amostras, fornecida pelo o servidor. \textcolor{red}{h  extra "espa os"...verifique isso....acho que vc poderia quebrar esse para em dois e explicar um pouco melhor}


\begin{figure}[H]
\begin{center}
			\caption{Diagrama Esquem tico do Sistema Proposto}
			\includegraphics[width=.8\textwidth]{Figuras/diaG.PNG}
            \vspace*{\fill} 
            \begin{quote} 
            \centering 
           Fonte: {Elaborada pela autora}
            \end{quote}
            \vspace*{\fill}
			\label{fig:dia}

\end{center}
\end{figure}



\subsection{DESENVOLVIMENTO DO HARDWARE}
\hspace*{0.8cm}Ap s a etapa de defini o e aquisi o de componentes que fazem parte da configura o do hardware, procedeu-se a implementa o do circuito eletr nico em um protoboard, sendo que inicialmente foram realizados testes nos componentes eletr nicos. O primeiro teste consistiu na verifica o dos sinais ECG utilizando a placa ECG AD8232, a qual foi alimentada com a tens o de 3,3V e utilizou o  pino GND para aterramento do sinal. Utilizou-se os pinos RA, LA e RL do simulador de sinais vitais HS-15 no modo ECG, conectado nos pinos RA, LA e RL da placa AD8232 com o intuito de analisar a forma de onda caracter stica do sinal ECG, com o manuseio da ponta de prova do oscilosc pio e assim analisar o sinal na tela do equipamento. \textcolor{red}{verifique os "espa os" novemente}


A segunda parte   realizar e estabelecer a comunica o entre o Raspberry Pi 3 e o MCP3008. Para a leitura de dados fornecedido pelo o conversor anal gico-digital, bibliotecas foram instaladas, utilizando os terminais do Raspberry (19, 21, 23, 24 e 26) para que ambos conectasse atrav s da conex o serial SPI que possuem. A conex o foi realizada da seguinte forma: os pinos VDD, VREF, AGND, CLK, DOUT, DIN, CS/SHDN e DGND do MCP3008 foram conectados ao Raspberry Pi de acordo  a pinagem estabalecida na Tabela 3. \textcolor{red}{corrija: fornecido, conectasse....a espa os extras novamente...}

\textcolor{red}{Os testes foram realizados em cada porta, ou seja, inserindo qualquer sinal as porta de entradas CH0, CH1, CH2, CH3, CH4, CH5, CH6, CH7 e os valores gerados no Raspberry em tempo real, qualquer altera o que   feito nas portas de entrada   imediatamente alterado a leitura da porta que esta sendo utilizada}. Os valores das 08 portas de entradas s o exibidos em uma tabela contendo os valores de canal e, a cada meio segundo, uma nova linha   impressa com os novos valores do canal. Cada coluna representa um canal espec fico da porta de entrada no MCP3008 e esses valores variam de 0 a 1023, onde 0 significa que o sinal est  no n vel do terra (GND) e 1023 significa que ele est  no valor VREF (3,3V) ou superior. \textcolor{red}{O TEXTO QUE EU MARQUEI EM VERMELHO DEVER  SER ESCRITO DE UM MODO MELHOR...parece que falta virgula e uma preposi o}

\begin{table}[H]
\centering
\caption{Conex o dos Pinos}
\begin{center}
\label{my-label}
\begin{tabular}{|c|c|}
\hline
\multicolumn{2}{|c|}{\textbf{Pinos}} \\ \hline
\multicolumn{1}{|l|}{\textbf{MCP3008}} & \multicolumn{1}{l|}{\textbf{Raspberry Pi}} \\ \hline
16 & 01 \\ \hline
15 & 01 \\ \hline
14 & 05 \\ \hline
13 & 23 \\ \hline
12 & 21 \\ \hline
11 & 19 \\ \hline
10 & 24 \\ \hline
9 & 39 \\ \hline
\end{tabular}
\end{center}
 \vspace*{\fill} 
            \begin{quote}
            \centering 
            \vspace{12pt}Fonte: {Elaborada pela autora} 
            \end{quote}
            \vspace*{\fill}
			\label{fig:ramcor}
\end{table}


A Figura \ref{fig:lop} representa a primeira disposi o de componentes no protoboard que foi desenvolvido atrav s de uma plataforma de software aberta chamada Fritzing \textcolor{red}{aqui vai uma ref} , na qual   possivel ao usu rio documentar prot tipo, esquematizar placa, layout de placa, entre outras fun es. 
\begin{figure}[H]
\begin{center}
			\caption{Circuito Esquem tico}
			\includegraphics[width=.6\textwidth]{Figuras/PROT.PNG}
			 \vspace*{\fill} 
            \begin{quote} 
            \centering 
           Fonte: {Elaborada pela autora}
            \end{quote}
            \vspace*{\fill}
			\label{fig:lop}
\end{center}
\end{figure}
O circuito esquem tico desenvolvido cont m os seguintes componentes: placa front-end \textcolor{red}{explique o que   isso}, MCP3008 e Raspeberry. A confec o da parte inferior da placa est  demonstrada especificamente na Figura \ref{fig:esque}. Nota-se que a placa possui dois segmentos tra ados que s o jumper soldado na parte superior da placa universal.
\begin{figure}[H]
\begin{center}
			\caption{Representa o inferior da trilha da placa universal}
			\includegraphics[width=.5\textwidth]{Figuras/ofic.png}
               \vspace*{\fill} 
            \begin{quote} 
            \centering 
           Fonte: Elaborada pela autora

            \end{quote}
            \vspace*{\fill}
			\label{fig:esque}
\end{center}
\end{figure}

Na pr tica esse circuito esquem tizado   montado no protoboard, conforme a Figura \ref{fig:esquem}, por m nessa imagem o conversor e o raspberry j  est o configurados para amostragem do sinal ECG. Destaca-se o excesso e a desordem de jumper conectados nos componentes e protoboard. Os sinais obtidos do simulador s o amplificados e tratados atrav s da placa fornecida pela  Sparkfun. Os dados s o lidos na porta e convertido em sinais digitais e transferidos para o Raspberry atrav s do protocolo SPI. \textcolor{red}{vc escreveu "esquematizado" de maneira errada}

Para a amplifica o dos sinais caracter stico a onda de ECG os pinos da Placa Front-End RA, LA e RL foram conectados especificamente em cada pino RA, LA e RL do simulador de sinais vitais no modo ECG. A conex o entre o HS-15 e a Placa Front-End   feita atrav s de cabos. Esses cabos s o de grande import ncia, pois isolam o paciente do circuito externo. Utilizou-se o produto cabo paciente ecg 5 vias da empresa Dixtal Philips, o qual foi modificado para o circuito, pois fornece apenas tr s vias RA, LA, RL e com isso retirou-se as vias V e LL. \textcolor{red}{h  erro de concord ncia em caracter stico}

Ap s a realiza o da montagem dos componentes no protoboard, observou-se a necessidade de eliminar jumper, pois causava interfer ncia no circuito. Assim, foi confeccionada uma placa de fabrica o manual com apenas  conectores: o primeiro conector tipo macho de 20 pinos para inserir o cabo flat do raspberry, conector f mea para inserir o pr prio MCP3008, 2 conectores para a placa front-end um conector f mea para os pinos GND, 3,3V e OUTPUT e outro macho para conectar com os pinos RA, LA e RL com o conector do cabo paciente.

\begin{figure}[H]
\begin{center}
			\caption{Representa o inferior da trilha da placa universal}
			\includegraphics[width=.5\textwidth]{Figuras/pop.jpeg}
             \vspace*{\fill} 
            \begin{quote} 
            \centering 
           Fonte: Elaborada pela autora
            \end{quote}
            \vspace*{\fill}
			\label{fig:esquem}
\end{center}
\end{figure}

Com base na representa o inferior da trilha da placa inferior foi poss vel desenvolver a mesma representa o na placa de circuito confeccionada para o projeto, para isso utilizou-se apenas solda de estanho-chumbo  e esta o de solda. A Figura \ref{fig:trilhainfesup} exibe a parte inferior e superior da placa confeccionada. Na esquerda tem-se a representa o inferior enquanto na direita tem-se a parte superior. \textcolor{red}{erro: espa os extras}


A regi o inferior tem-se apenas pinos de conex o e solda estanho-chumbo e a superior tem-se apenas os conectores para inserir os componentes para conex o da placa, conversores e Raspeberry. A placa confeccionada foi inserido em uma caixa para proteger o circuito.

\begin{figure}[H]
\begin{center}
			\caption{Representa o inferior da trilha da placa universal}
			\includegraphics[width=.9\textwidth]{Figuras/pl.PNG}
             \vspace*{\fill} 
            \begin{quote} 
            \centering 
         Fonte: Elaborada pela autora
            \end{quote}
            \vspace*{\fill}
			\label{fig:trilhainfesup}
\end{center}
\end{figure}

\subsection{FIRMWARE}
\hspace*{0.8cm}O Raspberry foi programado para ler os dados fornecido pelo o MCP3008 durante a utiliza o do simulador conectado a placa amplificadora. Para tal finalidade   necess rio fazer uso de bibliotecas que precisam esta \textcolor{red}{<-ESSA PALAVRA EST  ERRADA} instalada no terminal do Raspberry. Ap s a conex o  da fia o ser conectada, instalar a biblioteca Adafruit MCP3008 Python fornecido pela a plataforma aberta GitHub \textcolor{red}{AQUI VAI UMA REF}. A partir da fonte do GitHub   necess rio que o Raspberry esteja conectado a rede e acessar o terminal com os seguintes comandos nessa ordem especifico \textcolor{red}{<- AJEITE ESSA FRASE} :
\begin{figure}[H]
\begin{center}
		
			\includegraphics[width=.8\textwidth]{Figuras/com.PNG}
			
\end{center}
\end{figure}

Ao baixar as bibliotecas atrav s da plataforma   necess rio modificar a pasta de acesso para Adafruit/Python/MCP3008, nesse repositorio   necess rio modificar simpleteste.py e alterar a configura o do Software de acordo com a pinagem entre o conversor e o raspberry estabelecida:
\begin{itemize}
\item MCP3008 VDD para Raspberry Pi 3,3V
\item MCP3008 VREF paraa Raspberry Pi 3,3V
\item MCP3008 AGND para Raspberry Pi GND
\item MCP3008 DGND para Raspberry Pi GND
\item MCP3008 CLK para Raspberry Pi SCLK
\item MCP3008 DOUT para Raspberry Pi MISO 
\item MCP3008 DIN para Raspberry Pi MOSI
\item MCP3008 CS/SHDN Raspberry Pi CE0
\end{itemize}

Ao executar o comando na porta do raspberry, o programa imprimir  uma tabela com todos os canais ADC e seus valores. A cada meio segundo, uma nova linha ser  impressa com os mais recentes valores de canal.

Cada coluna representa um canal diferente e o cabe alho da primeira linha mostra o n mero do canal (de 0 a 7, 8 canais no total). O valor para cada canal   o valor ADC para esse canal. Este   um n mero que varia de 0 a 1023, onde 0 significa que o sinal est  no n vel do solo e 1023 significa que ele est  no valor VREF (3,3V) ou superior. Os outros valores s o obtidos por rela o proporcional, por exemplo, um valor de 512   3,3 / 2 ou 1,65 volts. \textcolor{red}{NIVEL DO SOLO????? N O D  PARA UTILIZAR OUTRO TERMO???}

O MCP3008 possui 8 portas de leitura de sinais, para a configura o do projeto usou-se a porta 04 para a leitura da placa de aquisi o, essa leitura   mostrada na Figura \ref{fig:ciru}. Para configura o cliente-servidor foi utilizado a estrutura web Pyhton Tornado vers o 4.5.2 fornecido pela fonte github. O tornado pode ser utilizado para conex es abertas, pesquisas longas, WebSocket ou aplica es que exigem uma conex o longa para o usu rio. Neste caso   utilizado o protocolo de comunica o HTTP-(\textit{HyperText Transfer Protocol}) \textcolor{red}{AQUI VAI UMA REF}, o qual permite a estrutura de cliente-servidor, proposta neste trabalho. Uma requisi o   caracterizada por uma mensagem HTTP enviada de um cliente para um servidor e a resposta   caracterizada por uma mensagem HTTP enviada por um servidor a um cliente.

A comunica o   feita atrav s da porta 8888. Utilizou-se os c digos em Phython do cliente e servidor, fornecido no AP NDICE A, utilizando como meio de comunica o o protocolo definido. Para isto configurou a rede Wi-Fi para o meio de comunica o. \textcolor{red}{VC ESCREVEU ERRADO "CONFIGUROU"...LEIA NOVAMENTE A FRASE}  

\begin{figure}[H]
\begin{center}
			\caption{Leitura da porta do MCP3008}
			\includegraphics[width=0.6\textwidth]{Figuras/D.PNG}
              \vspace*{\fill} 
            \begin{quote} 
            \centering 
           Fonte: Elaborada pela autora
            \end{quote}
            \vspace*{\fill}
			\label{fig:ciru}
\end{center}
\end{figure}
Para o uso das sintaxes do c digo cliente e servidor,   necess rio utilizar a fonte do github e importar algumas bibliotecas para manipula o de dados. Foram importados as bibliotecas contidos na c digo cliente e servidor para a utiliza o do c digo.

O cliente faz requisi o e o servidor o responde na ordem da solicita o. A fun o urllib2 presente no c digo cliente    uma fun o do Python usado para obter URLs. Nesse caso a busca   realizada usando protocolo 
HTTP, o qual   baseado em solicitar e responder. O acesso   feito pelo o endere o do IP do servidor acessando a porta 8888 do Raspberry pi 3. \textcolor{red}{ESPA OS NOVAMENTE KKKKK}

O diagrama da Figura \ref{fig:kenny}, baseia-se na seguinte proposta: o cliente faz requisi o ao servidor o qual acessa a porta 8888 do Raspberry. Nesta porta est  sendo executado o Tornado, o qual manipula a solicita o HTTP. O servidor comunica-se atrav s do protocolo SPI com o convers o MCP3008 e faz a leitura de 10000 pontos da porta 4 do MCP3008 fornecido pelo o sistema de aquisi o. Esses dados s o transmitidos ao cliente em forma de lista e o mesmo manipula para a gera o do gr fico.

 
 \begin{figure}[H]
\begin{center}
			\caption{Diagrama do projeto Cliente-Servidor}
			\includegraphics[width=1\textwidth]{blocos.PNG}
              \vspace*{\fill} 
            \begin{quote} 
            \centering 
           Fonte: {Elaborada pela autora} 
            \end{quote}
            \vspace*{\fill}
			\label{fig:lopioupe}
            \vspace*{\fill}
			\label{fig:kenny}
\end{center}
\end{figure}
 
\newpage

\section{RESULTADOS}

\hspace*{0.8cm} Neste cap tulo s o abordados os resultados obtidos durante o desenvolvimento do trabalho de conclus o de curso. Destaca-se neste projeto dois resultados significativos, um deles   o prot tipo f sico desenvolvido, apresentado na Figura \ref{fig:mama} e o outro resultado diz respeito as formas de ondas obtidas que comp e mo sinal de ECG, formando pela onda P, complexo QRS e onda T. \textcolor{red}{E O BAIXO CUSTO????}

Os componentes utilizados na montagem do prot tipo s o apresentados no cap tulo 3. Ap s, as etapas de montagem, soldagem e configura o, realizou o acomodamento do sistema, para isso utilizou uma caixa patola deixando apenas um espa o destinado conex o do cabo-flat com os pinos do Raspberry e a entrada do cabo paiciente, destaca-se que esse tipo de caixa   amplamente utilizada para acabamento est tico, prote o de circuitos. A Figura \ref{fig:mama}, representa o dispositivo final, disposto na caixa patola.




\begin{figure}[H]
\begin{center}
			\caption{Dispositivo final, disposto na caixa patola}
			\includegraphics[width=.9\textwidth]{Figuras/IMG_3651.jpg}
              \vspace*{\fill} 
            \begin{quote} 
            \centering 
           Fonte: Elaborada pela autora
            \end{quote}
            \vspace*{\fill}
			
			\label{fig:mama}
\end{center}
\end{figure}
A utiliza o da placa baseada no CI AD8232 da Analog Device para obten o dos sinais de ECG foi importante para eliminar as vari veis do processo de aquisi o. O prot tipo desenvolvido foi testado com o simulador de sinais vitais HS-15, e seus resultados foram comparados aos resultados da forma de onda padr o da segunda deriva o eletrocardiogr fica de um equipamento comercial modelo EP-12 da empresa Philips Dixtal \textcolor{red}{AQUI VAI UMA REF}.

A Figura \ref{fig:popi}, representa o sinal ECG simulado no equipamento eletrocardi grafo EP12, o equipamente utiliza o mesmo modelo de simulador de sinais vitais utilizando neste trabalho.   possivel, analisar e verificar a semelhan a com a o sinal exposto na Figura  \ref{fig:oscilo}.  O prop sito deste projeto   fornecer com qualidade e par metros as formas de ondas caracter stica do ECG. \textcolor{red}{MELHORE ESSA ULTIMA FRASE}

\begin{figure}[H]
\begin{center}
			\caption{Sinal de ECG simulado no EP12 Dixtal}
			\includegraphics[width=.8\textwidth]{Figuras/exame.PNG}
              \vspace*{\fill} 
            \begin{quote} 
            \centering 
           Fonte: \cite{ronald}
            \end{quote}
            \vspace*{\fill}
			\label{fig:popi}
\end{center}
\end{figure}

A Figura \ref{fig:oscilo}, representa a forma de onda adquirida pelo o sistema de aquisi o, utilizando o simulador para gerar sinais vitais e o oscil scopio para an lise. Este sinal   composto pela a onda p, complexo QRS e a onda T, caracter stico da forma de onda do eletrocardiograma. \textcolor{red}{MELHORE ESSE PARA...ESCREVA MAIS OU MENOS ASSIM: NA FIGURA TAL, APRESENTA-SE A FORMA DE ONDA....SACOU??...EM GERAL, PROCURE ESCREVER ASSIM...OU SEJA, CORRIJA O RESTO...}

\begin{figure}[H]
\begin{center}
			\caption{Sinal do ECG obtido pelo o oscilosc pio}
			\includegraphics[width=.9\textwidth]{Figuras/oscil.PNG}
             \vspace*{\fill} 
            \begin{quote} 
            \centering 
           Fonte: Elaborada pela autora
            \end{quote}
            \vspace*{\fill}
			\label{fig:oscilo}
\end{center}
\end{figure}

\hspace*{0.8cm}O Raspberry Pi 3 est  sendo usado como servidor. O cliente solicita atrav s do comando 'req' e para cada amostra   determinado um valor de acordo com o sinal fornecido pelo o simulador atrav s da porta do mcp.read-adc(4), ou seja, a porta 4 do convers o analogico/digital. Essa leitura est  sendo comprovado na Figura \ref{fig:kiki}.

\begin{figure}[H]
\begin{center}
			\caption{Amostragem Sinal do ECG}
			\includegraphics[width=.9\textwidth]{Figuras/json.PNG}
			              						\vspace*{\fill} 
            \begin{quote} 
            \centering 
           Fonte: Elaborada pela autora
            \end{quote}
            \vspace*{\fill}
			\label{fig:kiki}
\end{center}
\end{figure}

Os testes de transmiss o e recep o do sinal utilizando a arquitetura cliente-servidor proposta foram realizados conectando os dispositivos atrav s da rede Wi-fi do HUB. Com o resultado dessa arquitetura, tem-se a forma de onda utilizando o simulador de sinais vitais, conectado ao prot tipo. Nestes testes, utilizou-se o simulador e foi variado a amplitude e o batimento por minuto.

A Figura \ref{fig:ecg2}, representa um resultado da arquitetura e o prot tipo. A imagem representa o sinal ECG simulado, composto por 500 amostras de amplitude do sinal em determinado tempo. Esse sinal   composto inicialmente pela a onda P, complexo QRS e seguindo a onda T. 


\begin{figure}[H]
\begin{center}
			\caption{Sinal ECG simulado}
			\includegraphics[width=.8\textwidth]{Figuras/sinal.PNG}
                          						\vspace*{\fill} 
            \begin{quote} 
            \centering 
           Fonte: Elaborada pela autora
            \end{quote}
            \vspace*{\fill}
			\label{fig:ecg2}
\end{center}
\end{figure}
\begin{figure}[H]

\hspace*{0.8cm}Para a valida o do sistema proposto, o usu rio solicita atrav s de comando utilizando o cliente, o qual requisita ao servidor usando topologia estabelicida, o qual faz a leitura de dados e a transfere para o servidor. A Figura \ref{fig:syst}, exibe o sistema em opera o: uma tela est  exibindo a resposta solicitada ao servidor, ou seja, o sinal ECG adquirido pelo sistema de aquisi o utilizando sinais fornecido pelo o  simulador de sinais vitais e a outra tela, apenas faz a exibi o da leitura do MCP3008.

\begin{center}
			\caption{Sistema em Opera o}
			\includegraphics[width=.9\textwidth]{sist.PNG}
            	\vspace*{\fill} 
            \begin{quote} 
            \centering 
           Fonte: Elaborada pela autora
            \end{quote}
            \vspace*{\fill}
			\label{fig:syst}
\end{center}
\end{figure}

\newpage
\vspace*{4cm}
\begin{center}
\textbf{CONCLUS O}
\end{center}
\vspace{48pt}

\textcolor{red}{FALTOU FALAR DO BAIXO CUSTO...OU RETIRAMOS ESSA CONTRIBUI O...}


\addcontentsline{toc}{section}{CONCLUS O}

\hspace*{0.8cm}O objetivo deste trabalho foi desenvolver um  prot tipo de aquisi o de sinais eletrocardiogr ficos port til de baixo custo, para obter as ondas P, complexos QRS e ondas T. Inicialmente, realizou-se uma revis o das  reas de conhecimento relativo aos sinais que constituem o eletrocardiograma, estrutura de aquisi o e transmiss o de dados, tendo em vista atender a proposi o do trabalho de pesquisa.

O prot tipo desenvolvido, conseguiu atender as especifica es iniciais, gerando resultados satisfat rios que permitem a an lise do sinal e a comunica o com um computador em tempo real, atrav s apenas de uma solicita o utilizando uma rede de acesso sem fio, empregando o protocolo 802.11.

Com base no prot tipo desenvolvido,   poss vel analisar diretamente a forma de onda P, complexo QRS e onda T, e verificar a influ ncia do batimento card aco sobre o sinal. Foram realizados testes com a frequ ncia m nima de 30 btm e frequ ncia m xima de 240 btm, para assim,  determinar a oxigena o do corpo e poss veis disfun es   sa de, ou at  mesmo casos cr ticos que podem levar a morte repentina. \textcolor{red}{ESPA OS NOVAMENTE}

Ao decorrer do trabalho surgiram novas implementa es que  foram desenvolvidas, por m sem  xito, sendo uma destas a implementa o do cabo paciente. O prot tipo foi desenvolvido para receber a conex o do cabo paciente (AE-0H002-0) , por m os resultados da forma de onda ECG, n o foram satisfat rios \textcolor{red}{TU MOSTROU ISSO NO CAP DE RESULTADOS??? SE N O, COLOQUE}. Outras ideias para melhoria do sistema n o foram desenvolvidas, devido ao espa o de tempo limitado e a complexidade envolvida \textcolor{red}{<- TIRE ESSA FRASE}. Este trabalho permite diversas novas aplica es, e deixa em aberto possibilidades para melhorias e  desenvolvimento de outros trabalhos correlatos.

Com base nisso merecem destaque os seguintes trabalhos futuros.

\begin{itemize}
\item Eliminar os ru dos presente na forma de onda que comp e o ECG.
\item Melhoria do cabo paciente (AE-0H002-0).
\item Cria o de uma interface gr fica que atendam as necessidade do usu rio.
\end{itemize}

Independentemente destas etapas n o terem sido incorporadas no prot tipo, o mesmo j  foi projetado visando incorporar estas futuras melhorias, permitindo a continuidade do projeto.

\newpage


% Uncomment the following two lines if you want to have a bibliography. Please do not forget to add an entry to your bibliography and reference it by using the \cite{} command
%
\providecommand*{\refname}{}
\renewcommand*{\refname}{\begin{center}\textbf{REFER NCIAS}\end{center}}

\newpage
\vspace*{4cm}
%\providecommand{\bibname}{}
%\renewcommand{\bibname}{\begin{center}\textbf{REFER NCIAS}\vspace{48pt}\end{center}}

\pagestyle{empty}
\pagestyle{myheadings}

%

\addcontentsline{toc}{section}{REFER\^ENCIAS}


\newpage
%\pagestyle{fancy}
%\renewcommand{\headrulewidth}{0pt}
%\rhead{\thepage}
%\fancyfoot[C]{}
\begin{center}
\textbf{AP NDICE A-  CLIENTE/SERVIDOR}

\addcontentsline{toc}{section}{AP NDICE A - CLIENTE/SERVIDOR}


\textbf{CLIENTE}

\lstinputlisting[language={Python},inputencoding={utf8},extendedchars=false]{cliente.py}


\newpage
\textbf{SERVIDOR}

\lstinputlisting[language={Python},inputencoding={utf8},extendedchars=false]{servidor.py}

\end{center}

\newpage
\begin{center}
\textbf{AP NDICE B - SINAIS ECG OBTIDOS PELO O SISTEMA DE AQUISI O}

\addcontentsline{toc}{section}{AP NDICE B - SINAIS ECG OBTIDOS PELO O SISTEMA DE\newline AQUISI O}
\begin{itemize}
\item Amplitude 0,5mV
\begin{figure}[H]
\begin{center}
\caption{Sinal ECG 30 bpm}
					\includegraphics[width=0.7\textwidth]{Figure_1-1.png}
			     \vspace*{\fill} 
            \begin{quote} 
            \centering 
            Fonte: (Acervo da Autora, 2017)
            \end{quote}
            \vspace*{\fill}
			\label{fig:eletrodo}
\end{center}
\end{figure}



\newpage

\item Amplitude 1,5mV

\begin{figure}[H]
\begin{center}
\caption{Sinal ECG 30 bpm}
		
			\includegraphics[width=0.7\textwidth]{Figure_3-1.png}
       \vspace*{\fill} 
            \begin{quote} 
            \centering 
            Fonte: Elaborada pela autora
            \end{quote}
            \vspace*{\fill}
			\label{fig:eletrodo}
\end{center}
\end{figure}

\end{itemize}

\end{center}

% End of the document
\end{document}